%% toc: P15 174 174.9,174.10
% Structural-model form of the finite assignment experiment; identification and its boundaries.
\chapter{Randomized Assignment as a Structural Causal Model}
\label{sec:p15-w02b-p15-077-assignment-scm}

\Cref{sec:p15-factorial-causal} separates the effect of a randomized
assignment from a claim about the mechanism that produces it. This chapter
makes the first half of that separation exact in the vocabulary of finite
structural causal models. It writes the finite assignment experiment of
\cref{def:p15-factorial-assignment} as a model in the sense of
\cref{def:p16-finite-scm}, proves that the assignment estimand equals a
combination of intervention means of that model, and proves that it is
identified, in the sense of \cref{def:p16-causal-identification}, from the
joint law of the assignment and the recorded responses. It then proves where
this identification fails: when units interfere with one another, when the
effect of the treatment actually received is the target, and when outcome
recording depends on the response. The last section states which premises an
empirical reasoning study, such as the design of
\cref{def:p19-proposed-factorial-design}, must supply before these results
apply to it.

\paragraph{Inputs and conclusions.}
The inputs are a finite set of recruited units, a finite set of treatment
cells, a known assignment law with positive marginal cell probabilities, a
response for each unit at each cell, and fixed contrast weights. The
conclusions are equalities between the assignment estimand, intervention
means in the structural model, and an explicit functional of the observed
law; nonidentification statements; and a list of premises that an empirical
study must document. All probabilities are finite sums, in the sense of
\cref{sec:p15-finite-conditioning,sec:p15-finite-independence}.

\paragraph{Sources.}
The two-model criterion for identification is Pearl's Definition~2
\cite[sec.~3.2.1, p.~109]{Pearl2009CausalOverview}, adopted in
\cref{def:p16-causal-identification}. The argument that an exogenous
assignment independent of the response disturbance makes the intervention
mean equal to the conditional mean follows the derivation of
\cite[sec.~3.2.2, eqs.~(11)--(13), p.~109]{Pearl2009CausalOverview}.
Pearl models a trial with imperfect compliance by a randomized assignment,
a received treatment and a response, under an exclusion assumption and an
independence assumption secured by randomization, and records that the
effect of received treatment is not identified without additional
assumptions \cite[sec.~3.5.2, pp.~122--124]{Pearl2009CausalOverview}.
The finite multi-unit models, the identifying functional, the interference
and missingness results and all constructions below are stated and proved in
this handbook.

\section{The structural model of an assignment experiment}
\label{sec:p15-w02b-p15-077-model}

Throughout, \(N\geq1\) recruited units and a nonempty finite set
\(\mathcal T\) of treatment cells are fixed. An assignment vector is
\(z=(z_1,\ldots,z_N)\in\mathcal T^N\); for \(t\in\mathcal T\),
\(t\mathbf 1\) denotes the vector \((t,\ldots,t)\) assigning every unit to
cell \(t\). Vertex names below are labels; any enumeration identifies a
vertex set with the set \(\{1,\ldots,n\}\) of \cref{def:p16-causal-dag}.

\begin{hbdefinition}{def:p15-w02b-p15-077-structural-experiment}{The structural model of a finite assignment experiment}
Let \(E\) be a finite assignment experiment as in
\cref{def:p15-factorial-assignment}: units \(1,\ldots,N\), cells
\(\mathcal T\), real potential responses \(Y_i(t)\), and a randomized
assignment vector \(Z=(Z_1,\ldots,Z_N)\). Write
\(\varpi(z)=\Pr(Z=z)\) for the \emph{design law} on \(\mathcal T^N\),
so that
\[
 \pi_i(t)=\sum_{z:\,z_i=t}\varpi(z),
\]
and assume \(\pi_i(t)>0\) for every unit \(i\) and cell \(t\).
The \emph{structural model} \(M_E\) of \(E\) has the following data.
\begin{enumerate}
\item Vertices \(0,1,\ldots,N\) and edges \(0\to i\) for
      \(i=1,\ldots,N\). Vertex \(0\) carries the assignment vector and
      vertex \(i\) the recorded response of unit \(i\).
\item Value sets \(S_0=\mathcal T^N\) and
      \(S_i=\{Y_i(t):t\in\mathcal T\}\), a finite nonempty set of reals.
\item Outcome space \(\Omega=\mathcal T^N\) with mass
      \(q=\varpi\). Disturbances \(U_0(\omega)=\omega\), with
      \(\mathcal U_0=\mathcal T^N\), and constant disturbances
      \(U_i(\omega)=0\), with \(\mathcal U_i=\{0\}\).
\item Structural functions \(f_0(u)=u\) and
      \(f_i(z,0)=Y_i(z_i)\).
\end{enumerate}
Write \(Z=V_0\) and \(Y_i^{\mathrm{rec}}=V_i\) for the solution
coordinates. The randomized assignment is thus an exogenous parent of every
response, and its law \(\varpi\) is known because the experimenter fixes it.
The premises of \cref{def:p15-factorial-assignment} appear as properties of
the equations: \emph{no interference} is the requirement that \(f_i\)
reads only the coordinate \(z_i\); \emph{stable treatment versions} is
the requirement that each cell yields one value \(Y_i(t)\); and
\emph{consistency} is the declaration that the recorded response is
\(f_i\) evaluated at the assigned vector.
\index{assignment experiment!structural model}\index{randomization!as exogenous parent}
\end{hbdefinition}

For the response function
\(h(y_1,\ldots,y_N)=N^{-1}\sum_{i=1}^Ny_i\) and a model \(M\) on
these vertices, put
\[
 \mu_M(t)\defeq\mathbb E\bigl[h\bigl(V_1^{0=t\mathbf 1},\ldots,
   V_N^{0=t\mathbf 1}\bigr)\bigr],
\]
the mean response when every unit's assignment is set to \(t\) by the
intervention \(\operatorname{do}(V_0=t\mathbf 1)\) of
\cref{def:p16-causal-intervention}. For weights \(w=(w_t)_{t\in\mathcal T}\)
put \(\tau(M)=\sum_tw_t\mu_M(t)\).

\begin{hbtheorem}{thm:p15-w02b-p15-077-correspondence}{proved=annex:p15-w02b-p15-077-corr}{The assignment estimand is an intervention contrast}
Let \(E\) be a finite assignment experiment with design law \(\varpi\)
and let \(M_E\) be its structural model.
\begin{enumerate}
\item At every \(\omega\in\Omega\), the unique solution is
      \(V_0(\omega)=\omega\) and \(V_i(\omega)=Y_i(\omega_i)\). The
      observational law is
      \[
       p_{M_E}(z,y)=\begin{cases}\varpi(z),&y_i=Y_i(z_i)\text{ for every }i,\\
                    0,&\text{otherwise},\end{cases}
      \]
      which is the joint law of the assignment vector and the recorded
      responses \(Y_i(Z_i)\) of \(E\).
\item For every \(z\in\mathcal T^N\) and every \(\omega\),
      \(V_i^{0=z}(\omega)=Y_i(z_i)\). In particular
      \(V_i^{0=t\mathbf 1}=Y_i(t)\) at every outcome.
\item \(\mu_{M_E}(t)=N^{-1}\sum_iY_i(t)\), and the estimand of
      \cref{def:p15-factorial-assignment} satisfies
      \(\tau=\tau(M_E)\). If \(\sum_tw_t=0\), then for every
      \(t_0\in\mathcal T\),
      \[
       \tau=\sum_{t\ne t_0}w_t\bigl(\mu_{M_E}(t)-\mu_{M_E}(t_0)\bigr),
      \]
      and each bracket is the average intervention effect of
      \cref{def:p16-causal-intervention} for the response function \(h\)
      and the fixed interventions \(\operatorname{do}(V_0=t\mathbf 1)\)
      and \(\operatorname{do}(V_0=t_0\mathbf 1)\). For weights equal to
      one at \(t\), minus one at \(s\ne t\) and zero elsewhere,
      \(\tau\) is that single average intervention effect.
\item Conversely, let \(M\) be a finite structural model on the same graph,
      with \(S_0=\mathcal T^N\), real finite \(S_i\),
      \(f_0(u)=u\), \(\mathbb P(U_0=z)=\varpi(z)\), and response functions
      \(f_i(z,u)=g_i(z_i)\) for functions
      \(g_i:\mathcal T\to S_i\). With \(Y_i(t)\defeq g_i(t)\), the model
      \(M\) and the structural model of the resulting experiment assign the
      same probability to every \((z,y)\in\mathcal T^N\times\mathbb R^N\)
      observationally and under every intervention
      \(\operatorname{do}(V_0=z)\).
\end{enumerate}
\end{hbtheorem}

The intervention \(\operatorname{do}(V_0=t\mathbf 1)\) is a hypothetical
experiment in which all recruited units receive the same cell. Under a
completely randomized design with two or more nonempty cells, the design law
gives that configuration probability zero. The correspondence nevertheless
equates its mean with a quantity defined from the potential responses,
because no interference makes each unit's response depend only on its own
coordinate.

\section{Identification from the assigned-arm law}
\label{sec:p15-w02b-p15-077-identification}

\begin{hbdefinition}{def:p15-w02b-p15-077-classes}{Model classes for a randomized assignment}
Fix \(N\), \(\mathcal T\), a design law \(\varpi\) on
\(\mathcal T^N\) with \(\pi_i(t)>0\) for all \(i,t\), and weights
\(w\). The class \(\mathcal C_{\mathrm{int}}(\varpi)\) contains every
finite structural model \(M\) of \cref{def:p16-finite-scm} with
vertices \(0,1,\ldots,N\), edges \(0\to i\), \(S_0=\mathcal T^N\),
finite real \(S_i\), and the following two properties.
\begin{enumerate}
\item \emph{Known exogenous assignment.} \(\mathcal U_0=\mathcal T^N\),
      \(f_0(u)=u\), and \(\mathbb P(U_0=z)=\varpi(z)\) for every
      \(z\).
\item \emph{Randomization.} \(U_0\) is independent of the tuple
      \(\bar U=(U_1,\ldots,U_N)\):
      \(\mathbb P(U_0=z,\bar U=u)=\varpi(z)\mathbb P(\bar U=u)\)
      for all \(z,u\). The response disturbances may depend on one another.
\end{enumerate}
The subclass \(\mathcal C_{\mathrm{NI}}(\varpi)\) adds a third property.
\begin{enumerate}\setcounter{enumi}{2}
\item \emph{No interference, stable versions and consistency.} For each
      \(i\) there is \(g_i:\mathcal T\times\mathcal U_i\to S_i\) with
      \(f_i(z,u)=g_i(z_i,u)\).
\end{enumerate}
In both classes all vertices are observed, and the quantity of interest is
\(\tau(M)=\sum_tw_t\mu_M(t)\). In \(\mathcal C_{\mathrm{NI}}(\varpi)\)
the response of unit \(i\) at cell \(t\) is the random variable
\(g_i(t,U_i)\); the finite-population experiment is the case in which each
\(g_i\) ignores its disturbance.
\index{causal identification!randomized assignment}
\end{hbdefinition}

The structural model \(M_E\) of an experiment with design law \(\varpi\)
belongs to \(\mathcal C_{\mathrm{NI}}(\varpi)\): constant disturbances are
independent of every random variable. For a mass function \(p\) on
\(\mathcal T^N\times\mathbb R^N\) with finite support, define
\[
 \Phi_t(p)\defeq\frac1N\sum_{i=1}^N\frac1{\pi_i(t)}
   \sum_{(z,y)}p(z,y)\,\mathbf 1_{\{z_i=t\}}\,y_i,
 \qquad
 \Phi_w(p)\defeq\sum_{t\in\mathcal T}w_t\Phi_t(p).
\]
Both maps use only \(p\) and the known design probabilities.

\begin{hbtheorem}{thm:p15-w02b-p15-077-identification}{proved=annex:p15-w02b-p15-077-ident}{Identification of the assignment contrast}
Let \(M\in\mathcal C_{\mathrm{NI}}(\varpi)\) and write \(p_M\) for its
observational law on all vertices.
\begin{enumerate}
\item For every \(t\in\mathcal T\),
      \(\mu_M(t)=\Phi_t(p_M)\). Hence \(\tau(M)=\Phi_w(p_M)\), and
      \(\tau\) is identified from the observational law within
      \(\mathcal C_{\mathrm{NI}}(\varpi)\).
\item Let \(\delta_v\) denote point mass at \(v\). At every outcome,
      \(\Phi_w(\delta_{V(\omega)})\) equals
      \[
       \widehat\tau(\omega)=\frac1N\sum_i\sum_t
        \frac{w_t\mathbf 1_{\{Z_i(\omega)=t\}}V_i(\omega)}{\pi_i(t)},
      \]
      and \(\mathbb E[\widehat\tau]=\Phi_w(p_M)=\tau(M)\). For
      \(M=M_E\), \(V_i=Y_i(Z_i)\), so this is the estimator of
      \cref{thm:p15-factorial-assignment}.
\item Suppose the design is completely randomized with fixed positive cell
      sizes \(n_t\), \(\sum_tn_t=N\): \(\varpi\) is uniform on the
      vectors with exactly \(n_t\) coordinates equal to \(t\) for every
      \(t\). Then \(\pi_i(t)=n_t/N\), and with the assigned-arm mean
      \(\overline Y_t=n_t^{-1}\sum_{i:Z_i=t}V_i\),
      \[
       \mu_M(t)=\mathbb E\bigl[\overline Y_t\bigr],\qquad
       \tau(M)=\sum_tw_t\,\mathbb E\bigl[\overline Y_t\bigr].
      \]
\end{enumerate}
\end{hbtheorem}

\Cref{thm:p16-causal-positive-identification} does not apply here: it
requires positive mass at every joint value of all vertices, whereas
\(p_M\) vanishes at every pair \((z,y)\) inconsistent with the response
functions, and a completely randomized design gives mass zero to most
assignment vectors. The identification above uses instead the explicit
randomization and no-interference properties of the class. It treats the
observational law as known exactly. Part~2 records how the identifying
functional, applied to the single realized assignment and its recorded
responses, yields an unbiased estimator; its uncertainty is the
randomization variance displayed in \cref{thm:p15-factorial-assignment}.

\section{Where the identification fails}
\label{sec:p15-w02b-p15-077-boundaries}

Each of the three results below removes or changes one premise of
\(\mathcal C_{\mathrm{NI}}(\varpi)\) and exhibits the consequence through
the two-model criterion of \cref{def:p16-causal-identification}.

\begin{hbtheorem}{thm:p15-w02b-p15-077-interference}{proved=annex:p15-w02b-p15-077-interf}{Interference: identification exactly when the common assignments have positive mass}
Work in \(\mathcal C_{\mathrm{int}}(\varpi)\), where a response function
may read the whole assignment vector.
\begin{enumerate}
\item If \(\varpi(t\mathbf 1)>0\), then for every model in the class
      \[
       \mu_M(t)=\mathbb E\bigl[h(V_1,\ldots,V_N)\mid V_0=t\mathbf 1\bigr].
      \]
\item If \(\varpi(t^*\mathbf 1)=0\), then for every
      \(M\in\mathcal C_{\mathrm{int}}(\varpi)\) and every real
      \(c\ne0\) there is \(M'\in\mathcal C_{\mathrm{int}}(\varpi)\)
      with \(p_{M'}=p_M\), \(\mu_{M'}(t^*)=\mu_M(t^*)+c/N\), and
      \(\mu_{M'}(t)=\mu_M(t)\) for every \(t\ne t^*\).
\item Consequently \(\tau\) is identified within
      \(\mathcal C_{\mathrm{int}}(\varpi)\) if and only if
      \(\varpi(t\mathbf 1)>0\) for every \(t\) with \(w_t\ne0\).
      Under complete randomization with at least two cells of positive
      size, \(\varpi(t\mathbf 1)=0\) for every \(t\), so no assignment
      estimand with a nonzero weight is identified in this class. In
      particular, for such a design and such weights, some
      \(M\in\mathcal C_{\mathrm{int}}(\varpi)\) has
      \(\Phi_w(p_M)\ne\tau(M)\), where \(\Phi_w\) is the functional
      of \cref{thm:p15-w02b-p15-077-identification}.
\end{enumerate}
\end{hbtheorem}

When interference is possible, the phrase ``the response of unit \(i\) at
cell \(t\)'' has no single referent, because the response may vary with
the other units' cells. The common-assignment mean \(\mu_M(t)\) remains a
well-defined intervention quantity, but a design that never assigns all
units together supplies no observation of it.

\begin{hbtheorem}{thm:p15-w02b-p15-077-noncompliance}{proved=annex:p15-w02b-p15-077-noncomp}{Noncompliance: the assignment effect is identified, the received-treatment effect is not}
Let \(\mathcal K(\varpi)\) be the class of finite structural models with
vertices \(0\), \(d_1,\ldots,d_N\), \(y_1,\ldots,y_N\), edges
\(0\to d_i\) and \(d_i\to y_i\), value sets \(S_0=\mathcal T^N\),
\(S_{d_i}=\mathcal T\) and finite real \(S_{y_i}\), with these properties:
the assignment equation and law are those of
\cref{def:p15-w02b-p15-077-classes}; \(U_0\) is independent of the tuple
of all other disturbances; the received treatment is
\(D_i=c_i(Z_i,U_{d_i})\) for some function \(c_i\); and the response is
\(Y_i=f_{y_i}(D_i,U_{y_i})\), so that assignment reaches the response only
through the received treatment. All vertices are observed. Write \(Y^{0=t\mathbf 1}\) for the tuple of
response coordinates under \(\operatorname{do}(V_0=t\mathbf 1)\), and
\(Y^{d=s\mathbf 1}\) for that tuple under the simultaneous intervention
setting every vertex \(d_i\) to \(s\). Put
\(\mu_M(t)=\mathbb E[h(Y^{0=t\mathbf 1})]\) and
\(\nu_M(s)=\mathbb E[h(Y^{d=s\mathbf 1})]\).
\begin{enumerate}
\item For every \(M\in\mathcal K(\varpi)\) and \(t\in\mathcal T\),
      \(\mu_M(t)=\Phi_t(p_{M,\{0,y_1,\ldots,y_N\}})\). The assignment
      contrast is identified.
\item Let \(\mathcal T=\{0,1\}\). There are \(M_A,M_B\in\mathcal K(\varpi)\)
      with the same observational law on all vertices, with
      \(\mu_{M_A}(1)-\mu_{M_A}(0)=\mu_{M_B}(1)-\mu_{M_B}(0)=0\), with
      \(\mathbb P(D_i=1)>0\) and \(\mathbb P(D_i=0)>0\) for every
      unit, but with
      \[
       \nu_{M_A}(1)-\nu_{M_A}(0)=0,\qquad
       \nu_{M_B}(1)-\nu_{M_B}(0)=\tfrac12.
      \]
      The received-treatment contrast is therefore not identified within
      \(\mathcal K(\varpi)\).
\end{enumerate}
\end{hbtheorem}

In the construction, the disturbance that governs uptake also enters the
response. This shared information is the dependence between compliance and
response factors that randomized assignment does not remove. Analyzing
by assigned cell estimates the assignment contrast; it is not the
received-treatment contrast unless further premises connect them.

\begin{hbtheorem}{thm:p15-w02b-p15-077-missingness}{proved=annex:p15-w02b-p15-077-miss}{Missing outcomes: response-independent loss and response-dependent loss}
Let \(\mathcal W(\varpi)\) be the class of finite structural models with
vertices \(0\), \(y_i\), \(r_i\), \(o_i\) for \(i=1,\ldots,N\), edges
\(0\to y_i\), \(0\to r_i\), \(y_i\to r_i\), \(y_i\to o_i\),
\(r_i\to o_i\), and the following equations. The assignment equation and
law are those of \cref{def:p15-w02b-p15-077-classes}, and \(U_0\) is
independent of the tuple of all other disturbances. The response is
\(Y_i=g_i(Z_i,U_{y_i})\) with real values. The recording indicator is
\(R_i=\rho_i(Z_i,Y_i,U_{r_i})\in\{0,1\}\). The recorded value is
\(O_i=Y_i\) if \(R_i=1\) and \(O_i=\star\) if \(R_i=0\), for a fixed
symbol \(\star\). The observed vertices are \(0\), the \(r_i\) and the
\(o_i\); the \(y_i\) are not observed. The quantity of interest is
\(\tau(M)=\sum_tw_t\mu_M(t)\) with
\(\mu_M(t)=\mathbb E[h(Y^{0=t\mathbf 1})]\).
\begin{enumerate}
\item Let \(\mathcal T=\{0,1\}\) and \(w=(w_0,w_1)=(-1,1)\). There are
      \(M_A,M_B\in\mathcal W(\varpi)\) with the same observational law on
      the observed vertices, the same recording probability in each arm,
      and \(\tau(M_A)=1\), \(\tau(M_B)=1/2\). In \(M_A\) recording
      depends on the assigned arm but not on the response; in \(M_B\) it
      depends on the response.
\item Let \(\mathcal W_0(\varpi)\) be the subclass in which
      \(\rho_i(t,y,u)\) does not depend on \(y\), and for every \(i\) the
      three disturbances \(U_0,U_{y_i},U_{r_i}\) are mutually independent.
      Assume \(\bar\rho_i(t)=\mathbb P(\rho_i(t,U_{r_i})=1)>0\) for all
      \(i,t\). Reading \(R_iO_i\) as zero when \(R_i=0\),
      \[
       \mu_M(t)=\frac1N\sum_{i=1}^N
        \frac{\mathbb E\bigl[\mathbf 1_{\{Z_i=t\}}R_iO_i\bigr]}
             {\mathbb P(Z_i=t,R_i=1)}
      \]
      for every \(M\in\mathcal W_0(\varpi)\), so \(\tau\) is identified
      within \(\mathcal W_0(\varpi)\). The model \(M_A\) of part~1 lies in
      this subclass.
\end{enumerate}
\end{hbtheorem}

Part~1 shows that loss confined to one arm cannot, from the observed law,
be distinguished from loss that depends on the response. Part~2 states a
sufficient premise under which loss, even when it differs between arms,
leaves the assignment contrast identified. That premise concerns the
unobserved responses and is not itself checked by the observed law.

\begin{hbexample}{def:p15-w02b-p15-077-worked}{Four participants, one per cell}
Take the cells \(\mathcal T=\{00,10,01,11\}\) of
\cref{def:p15-factorial-coding}, \(N=4\), and complete randomization with
\(n_t=1\): the design law gives mass \(1/24\) to each of the
\(24\) assignment vectors that place one participant in each cell, so
\(\pi_i(t)=1/4\). Stipulate the potential responses, in cell order
\((00,10,01,11)\),
\[
\begin{array}{c|cccc}
 \text{unit}&00&10&01&11\\\hline
 1&4&1&1&2\\
 2&3&0&2&2\\
 3&4&2&0&3\\
 4&3&1&1&1
\end{array}
\]
These are not observations from a conducted study. By
\cref{thm:p15-w02b-p15-077-correspondence},
\(\mu_{M_E}=(7/2,1,1,2)\). The aid-present checking weights
\((0,0,-1,1)\) give \(\tau=2-1=1\); the equal-weight weights
\((-1,1,-1,1)/2\) give \(\tau=((1-7/2)+(2-1))/2=-3/4\).

If unit \(1\) is assigned \(00\), unit \(2\) \(10\), unit \(3\)
\(01\) and unit \(4\) \(11\), the recorded responses are
\((4,0,0,1)\), and the aid-present estimate is \(1-0=1\). If instead
unit \(1\) is assigned \(11\) and unit \(2\) \(01\), with units
\(3,4\) in \(00,10\), the recorded responses give \(2-2=0\). Over all
\(24\) vectors, each unit occupies cell \(11\) and cell \(01\) in six
vectors each, so the mean estimate is
\((8/4)-(4/4)=1=\tau\), as \cref{thm:p15-w02b-p15-077-identification}
requires. In this design \(\overline Y_t\) is the single recorded response
in cell \(t\).

No assignment vector here puts all four units in cell \(11\). Changing
unit \(4\)'s response function only at that vector, from \(1\) to \(4\),
leaves the observational law unchanged and raises \(\mu(11)\) to
\(11/4\), so the aid-present contrast becomes \(7/4\). The model obtained
is an interference model; it is excluded by the no-interference premise and
illustrates part~2 of \cref{thm:p15-w02b-p15-077-interference}.
\end{hbexample}

\section{Premises an empirical reasoning study must supply}
\label{sec:p15-w02b-p15-077-studies}

\begin{hbmethod}{def:p15-w02b-p15-077-study-premises}{Applying the identification to an empirical reasoning study}
To use \cref{thm:p15-w02b-p15-077-identification} for a study such as
\cref{def:p19-proposed-factorial-design}, take the recruited participants
as the units, the four crossed conditions as \(\mathcal T\), and the
participant's scored primary outcome as the response. Declare the weights
before outcomes are examined: \((0,0,-1,1)\) for the aid-present checking
contrast and \((-1,1,-1,1)/2\) for the equal-weight contrast, in cell order
\((00,10,01,11)\). The study must then document the following premises,
each of which is a property of \(\mathcal C_{\mathrm{NI}}(\varpi)\) or of
\(\mathcal W_0(\varpi)\).
\begin{enumerate}
\item \emph{Assignment law.} The randomization procedure, its cell sizes
      or blocking, and hence \(\varpi\) and every \(\pi_i(t)\), fixed
      before assignment; the device is not informed by participant
      characteristics or anticipated responses, which secures the
      independence of \(U_0\) from the response disturbances; every
      \(\pi_i(t)\) is positive.
\item \emph{Stable versions.} Each condition is one fixed instruction text,
      aid presentation, timing and incentive. Variation not fixed by the
      protocol, such as experimenter or session, is either absent or made
      part of the cell definition.
\item \emph{Consistency.} The scored outcome is computed by the fixed
      scoring and tolerance rule from the responses produced under the
      assigned condition.
\item \emph{No interference.} Participants are tested without contact
      that could carry one participant's condition to another's responses.
      If participants share sessions, either the session is the unit of
      assignment or a model of the interference is supplied; with
      complete randomization of individuals,
      \cref{thm:p15-w02b-p15-077-interference} shows that the data alone
      cannot replace this premise.
\item \emph{Outcome recording.} Every primary outcome is recorded, or loss
      is reported by arm together with a premise of the form of
      \(\mathcal W_0(\varpi)\): recording does not depend on the
      response, given the assigned condition, and has positive probability.
      \Cref{thm:p15-w02b-p15-077-missingness} shows that this premise is
      not tested by the observed law.
\end{enumerate}
Under these premises the target is the assignment contrast in the recruited
sample. Three further claims need further premises. A received-treatment
claim needs premises about adherence beyond exclusion and randomization,
by \cref{thm:p15-w02b-p15-077-noncompliance}. A population claim needs a
sampling or transport argument. A mechanism claim needs the additional
evidence described after \cref{thm:p15-factorial-mechanism}.
\end{hbmethod}

\begin{hbcorollary}{thm:p15-w02b-p15-077-study-contrasts}{proved=annex:p15-w02b-p15-077-study}{The two declared contrasts of the factorial reasoning design}
Suppose the design of \cref{def:p19-proposed-factorial-design} assigns
participants by complete randomization with fixed positive cell sizes, and
its structural model lies in \(\mathcal C_{\mathrm{NI}}(\varpi)\).
Then the aid-present contrast \(\mu(11)-\mu(01)\) and the equal-weight
contrast \(\bigl((\mu(10)-\mu(00))+(\mu(11)-\mu(01))\bigr)/2\) are
identified, and they equal
\(\mathbb E[\overline Y_{11}-\overline Y_{01}]\) and
\(\mathbb E\bigl[\bigl((\overline Y_{10}-\overline Y_{00})
 +(\overline Y_{11}-\overline Y_{01})\bigr)/2\bigr]\), respectively.
The statistics \(\widehat\Delta_{\rm aid}\) and
\(\widehat\Delta_{\rm check}\) of that design are unbiased for them over
the randomization.
\end{hbcorollary}

\begin{hbwarning}{def:p15-w02b-p15-077-limits}{What the identification does not deliver}
The identified quantity is a contrast of common-assignment means for the
recruited units, under the stated equations. It is not a statement about
another population, about the effect of the treatment actually received, or
about the process through which assignment acts:
\cref{thm:p15-factorial-mechanism} gives two models with the same
assignment surface and different responses to an intervention on a
putative reflection variable. Identification also treats the observational
law as exactly known. An observable variance estimator, and an interval
with stated coverage, need the further design or model assumptions noted
after \cref{thm:p15-factorial-assignment}.
\end{hbwarning}

\paragraph{Relationships.}
\Cref{thm:p16-binary-po-bridge} relates potential outcomes and
interventions for one binary treatment, with a unit drawn from a joint law;
the present construction keeps the recruited units fixed, places the
randomness in the known assignment law, and allows any finite number of
cells. \Cref{thm:p16-po-ate-identification} identifies an average
treatment effect under conditional exchangeability; here the corresponding
independence is a property of the model class, secured by the
randomization premise, rather than an assumption about an observational
population. \Cref{thm:p16-causal-nonidentification} gives two-model
nonidentification for confounded treatment, and the three boundary results
above give the analogous constructions for interference, noncompliance and
response-dependent loss.
